% Estilo común del SRS, el SAD y el plan de pruebas. Reproduce el formato del SRS y el
% SAD originales: carta, márgenes de una pulgada, Times de 10 puntos en el cuerpo.
% Las rutas son relativas a la carpeta del documento: se compila desde ella
% (latexmk -cd) y se carga con % Estilo común del SRS, el SAD y el plan de pruebas. Reproduce el formato del SRS y el
% SAD originales: carta, márgenes de una pulgada, Times de 10 puntos en el cuerpo.
% Las rutas son relativas a la carpeta del documento: se compila desde ella
% (latexmk -cd) y se carga con % Estilo común del SRS, el SAD y el plan de pruebas. Reproduce el formato del SRS y el
% SAD originales: carta, márgenes de una pulgada, Times de 10 puntos en el cuerpo.
% Las rutas son relativas a la carpeta del documento: se compila desde ella
% (latexmk -cd) y se carga con % Estilo común del SRS, el SAD y el plan de pruebas. Reproduce el formato del SRS y el
% SAD originales: carta, márgenes de una pulgada, Times de 10 puntos en el cuerpo.
% Las rutas son relativas a la carpeta del documento: se compila desde ella
% (latexmk -cd) y se carga con \input{../comun/estilo}.

\usepackage[T1]{fontenc}
\usepackage[utf8]{inputenc}
\usepackage[spanish,es-tabla,es-noshorthands,es-lcroman]{babel}
\usepackage{tgtermes}
\usepackage{tgheros}
\usepackage[letterpaper,margin=1in]{geometry}
\usepackage{graphicx}
\usepackage{array}
\usepackage{tabularx}
\usepackage{setspace}
\usepackage{titlesec}
\usepackage{titletoc}
\usepackage{caption}
\usepackage[hidelinks]{hyperref}

\setstretch{1.15}
\setlength{\parindent}{0pt}
\setlength{\parskip}{6pt plus 2pt}

\addto\captionsspanish{%
  \renewcommand{\contentsname}{Tabla de Contenidos}%
  \renewcommand{\listtablename}{Lista de Tablas}%
  \renewcommand{\listfigurename}{Lista de Figuras}%
}

\titleformat{\section}{\normalfont\bfseries\fontsize{12}{14}\selectfont}{\thesection.}{0.3em}{}
\titleformat{\subsection}{\normalfont\bfseries\fontsize{11}{13}\selectfont}{\thesubsection}{0.3em}{}
\titleformat{\subsubsection}{\normalfont\bfseries\normalsize}{\thesubsubsection}{0.3em}{}
\titlespacing*{\section}{0pt}{14pt}{8pt}
\titlespacing*{\subsection}{0pt}{12pt}{6pt}
\titlespacing*{\subsubsection}{0pt}{10pt}{4pt}

\titlecontents{section}[0pt]{\addvspace{2pt}\bfseries}
  {\thecontentslabel.\ }{}{\titlerule*[0.3em]{.}\contentspage}
\titlecontents{subsection}[1.5em]{}
  {\thecontentslabel\ }{}{\titlerule*[0.3em]{.}\contentspage}
\titlecontents{subsubsection}[3em]{}
  {\thecontentslabel\ }{}{\titlerule*[0.3em]{.}\contentspage}
\titlecontents{table}[0.75in]{}
  {\tablename~\thecontentslabel.\ }{}{\titlerule*[0.3em]{.}\contentspage}
\titlecontents{figure}[0.75in]{}
  {\figurename~\thecontentslabel.\ }{}{\titlerule*[0.3em]{.}\contentspage}

\captionsetup{labelsep=period,font=it,justification=centering}

\newcommand{\kpNombre}{}
\newcommand{\kpVersion}{}
\newcommand{\kpFecha}{}
\newcommand{\NombreDocumento}[1]{\renewcommand{\kpNombre}{#1}}
\newcommand{\VersionDocumento}[1]{\renewcommand{\kpVersion}{#1}}
\newcommand{\FechaDocumento}[1]{\renewcommand{\kpFecha}{#1}}

\newcommand{\Portada}{%
  \hypersetup{
    pdftitle={\kpNombre\ - Kit Pagos Colombia},
    pdfauthor={Andres Henao Niño, Joan Emmanuel Orduz Chia, Joshua Prieto Zambrano, David Estevan Rodriguez Jurado}
  }%
  \begin{titlepage}
    \centering
    \setstretch{1}%
    \setlength{\parskip}{0pt}%
    \fontsize{12}{16}\selectfont
    \includegraphics[height=6.96cm]{../comun/logo-javeriana}\par
    \vspace{46pt}
    {\fontsize{14}{16}\selectfont\bfseries\kpNombre\par}
    \vspace{16pt}
    Presentan\par
    \vspace{16pt}
    Andres Henao Niño\par
    Joan Emmanuel Orduz Chia\par
    Joshua Prieto Zambrano\par
    David Estevan Rodriguez Jurado\par
    \vspace{31pt}
    Director\par
    Jaime Andrés Pavlich Mariscal\par
    \vspace{21pt}
    Pontificia Universidad Javeriana\par
    Facultad de Ingeniería\par
    Trabajo de Grado\par
    \vspace{16pt}
    Bogotá D.C.\par
    \kpFecha\par
    Versión \kpVersion\par
  \end{titlepage}%
  \pagenumbering{roman}%
}

% Cada fila se agrega con \Cambio{versión}{fecha}{descripción}{autor(es)}.
\newenvironment{HistorialDeCambios}{%
  \section*{Historial de Cambios}%
  \renewcommand{\arraystretch}{1.3}%
  \tabularx{\textwidth}{|c|c|X|>{\raggedright\arraybackslash}p{4.2cm}|}
  \hline
  \textbf{Versión} & \textbf{Fecha} & \textbf{Descripción del cambio} & \textbf{Autor(es)} \\
  \hline
}{%
  \endtabularx
  \clearpage
}
\newcommand{\Cambio}[4]{#1 & #2 & #3 & #4 \\ \hline}

% La plantilla PTP pide que la numeración arábiga empiece en la Introducción y que las
% listas de tablas y figuras no entren en la tabla de contenidos.
\newcommand{\Indices}{%
  \tableofcontents
  \clearpage
  \listoftables
  \listoffigures
  \clearpage
  \pagenumbering{arabic}%
}
.

\usepackage[T1]{fontenc}
\usepackage[utf8]{inputenc}
\usepackage[spanish,es-tabla,es-noshorthands,es-lcroman]{babel}
\usepackage{tgtermes}
\usepackage{tgheros}
\usepackage[letterpaper,margin=1in]{geometry}
\usepackage{graphicx}
\usepackage{array}
\usepackage{tabularx}
\usepackage{setspace}
\usepackage{titlesec}
\usepackage{titletoc}
\usepackage{caption}
\usepackage[hidelinks]{hyperref}

\setstretch{1.15}
\setlength{\parindent}{0pt}
\setlength{\parskip}{6pt plus 2pt}

\addto\captionsspanish{%
  \renewcommand{\contentsname}{Tabla de Contenidos}%
  \renewcommand{\listtablename}{Lista de Tablas}%
  \renewcommand{\listfigurename}{Lista de Figuras}%
}

\titleformat{\section}{\normalfont\bfseries\fontsize{12}{14}\selectfont}{\thesection.}{0.3em}{}
\titleformat{\subsection}{\normalfont\bfseries\fontsize{11}{13}\selectfont}{\thesubsection}{0.3em}{}
\titleformat{\subsubsection}{\normalfont\bfseries\normalsize}{\thesubsubsection}{0.3em}{}
\titlespacing*{\section}{0pt}{14pt}{8pt}
\titlespacing*{\subsection}{0pt}{12pt}{6pt}
\titlespacing*{\subsubsection}{0pt}{10pt}{4pt}

\titlecontents{section}[0pt]{\addvspace{2pt}\bfseries}
  {\thecontentslabel.\ }{}{\titlerule*[0.3em]{.}\contentspage}
\titlecontents{subsection}[1.5em]{}
  {\thecontentslabel\ }{}{\titlerule*[0.3em]{.}\contentspage}
\titlecontents{subsubsection}[3em]{}
  {\thecontentslabel\ }{}{\titlerule*[0.3em]{.}\contentspage}
\titlecontents{table}[0.75in]{}
  {\tablename~\thecontentslabel.\ }{}{\titlerule*[0.3em]{.}\contentspage}
\titlecontents{figure}[0.75in]{}
  {\figurename~\thecontentslabel.\ }{}{\titlerule*[0.3em]{.}\contentspage}

\captionsetup{labelsep=period,font=it,justification=centering}

\newcommand{\kpNombre}{}
\newcommand{\kpVersion}{}
\newcommand{\kpFecha}{}
\newcommand{\NombreDocumento}[1]{\renewcommand{\kpNombre}{#1}}
\newcommand{\VersionDocumento}[1]{\renewcommand{\kpVersion}{#1}}
\newcommand{\FechaDocumento}[1]{\renewcommand{\kpFecha}{#1}}

\newcommand{\Portada}{%
  \hypersetup{
    pdftitle={\kpNombre\ - Kit Pagos Colombia},
    pdfauthor={Andres Henao Niño, Joan Emmanuel Orduz Chia, Joshua Prieto Zambrano, David Estevan Rodriguez Jurado}
  }%
  \begin{titlepage}
    \centering
    \setstretch{1}%
    \setlength{\parskip}{0pt}%
    \fontsize{12}{16}\selectfont
    \includegraphics[height=6.96cm]{../comun/logo-javeriana}\par
    \vspace{46pt}
    {\fontsize{14}{16}\selectfont\bfseries\kpNombre\par}
    \vspace{16pt}
    Presentan\par
    \vspace{16pt}
    Andres Henao Niño\par
    Joan Emmanuel Orduz Chia\par
    Joshua Prieto Zambrano\par
    David Estevan Rodriguez Jurado\par
    \vspace{31pt}
    Director\par
    Jaime Andrés Pavlich Mariscal\par
    \vspace{21pt}
    Pontificia Universidad Javeriana\par
    Facultad de Ingeniería\par
    Trabajo de Grado\par
    \vspace{16pt}
    Bogotá D.C.\par
    \kpFecha\par
    Versión \kpVersion\par
  \end{titlepage}%
  \pagenumbering{roman}%
}

% Cada fila se agrega con \Cambio{versión}{fecha}{descripción}{autor(es)}.
\newenvironment{HistorialDeCambios}{%
  \section*{Historial de Cambios}%
  \renewcommand{\arraystretch}{1.3}%
  \tabularx{\textwidth}{|c|c|X|>{\raggedright\arraybackslash}p{4.2cm}|}
  \hline
  \textbf{Versión} & \textbf{Fecha} & \textbf{Descripción del cambio} & \textbf{Autor(es)} \\
  \hline
}{%
  \endtabularx
  \clearpage
}
\newcommand{\Cambio}[4]{#1 & #2 & #3 & #4 \\ \hline}

% La plantilla PTP pide que la numeración arábiga empiece en la Introducción y que las
% listas de tablas y figuras no entren en la tabla de contenidos.
\newcommand{\Indices}{%
  \tableofcontents
  \clearpage
  \listoftables
  \listoffigures
  \clearpage
  \pagenumbering{arabic}%
}
.

\usepackage[T1]{fontenc}
\usepackage[utf8]{inputenc}
\usepackage[spanish,es-tabla,es-noshorthands,es-lcroman]{babel}
\usepackage{tgtermes}
\usepackage{tgheros}
\usepackage[letterpaper,margin=1in]{geometry}
\usepackage{graphicx}
\usepackage{array}
\usepackage{tabularx}
\usepackage{setspace}
\usepackage{titlesec}
\usepackage{titletoc}
\usepackage{caption}
\usepackage[hidelinks]{hyperref}

\setstretch{1.15}
\setlength{\parindent}{0pt}
\setlength{\parskip}{6pt plus 2pt}

\addto\captionsspanish{%
  \renewcommand{\contentsname}{Tabla de Contenidos}%
  \renewcommand{\listtablename}{Lista de Tablas}%
  \renewcommand{\listfigurename}{Lista de Figuras}%
}

\titleformat{\section}{\normalfont\bfseries\fontsize{12}{14}\selectfont}{\thesection.}{0.3em}{}
\titleformat{\subsection}{\normalfont\bfseries\fontsize{11}{13}\selectfont}{\thesubsection}{0.3em}{}
\titleformat{\subsubsection}{\normalfont\bfseries\normalsize}{\thesubsubsection}{0.3em}{}
\titlespacing*{\section}{0pt}{14pt}{8pt}
\titlespacing*{\subsection}{0pt}{12pt}{6pt}
\titlespacing*{\subsubsection}{0pt}{10pt}{4pt}

\titlecontents{section}[0pt]{\addvspace{2pt}\bfseries}
  {\thecontentslabel.\ }{}{\titlerule*[0.3em]{.}\contentspage}
\titlecontents{subsection}[1.5em]{}
  {\thecontentslabel\ }{}{\titlerule*[0.3em]{.}\contentspage}
\titlecontents{subsubsection}[3em]{}
  {\thecontentslabel\ }{}{\titlerule*[0.3em]{.}\contentspage}
\titlecontents{table}[0.75in]{}
  {\tablename~\thecontentslabel.\ }{}{\titlerule*[0.3em]{.}\contentspage}
\titlecontents{figure}[0.75in]{}
  {\figurename~\thecontentslabel.\ }{}{\titlerule*[0.3em]{.}\contentspage}

\captionsetup{labelsep=period,font=it,justification=centering}

\newcommand{\kpNombre}{}
\newcommand{\kpVersion}{}
\newcommand{\kpFecha}{}
\newcommand{\NombreDocumento}[1]{\renewcommand{\kpNombre}{#1}}
\newcommand{\VersionDocumento}[1]{\renewcommand{\kpVersion}{#1}}
\newcommand{\FechaDocumento}[1]{\renewcommand{\kpFecha}{#1}}

\newcommand{\Portada}{%
  \hypersetup{
    pdftitle={\kpNombre\ - Kit Pagos Colombia},
    pdfauthor={Andres Henao Niño, Joan Emmanuel Orduz Chia, Joshua Prieto Zambrano, David Estevan Rodriguez Jurado}
  }%
  \begin{titlepage}
    \centering
    \setstretch{1}%
    \setlength{\parskip}{0pt}%
    \fontsize{12}{16}\selectfont
    \includegraphics[height=6.96cm]{../comun/logo-javeriana}\par
    \vspace{46pt}
    {\fontsize{14}{16}\selectfont\bfseries\kpNombre\par}
    \vspace{16pt}
    Presentan\par
    \vspace{16pt}
    Andres Henao Niño\par
    Joan Emmanuel Orduz Chia\par
    Joshua Prieto Zambrano\par
    David Estevan Rodriguez Jurado\par
    \vspace{31pt}
    Director\par
    Jaime Andrés Pavlich Mariscal\par
    \vspace{21pt}
    Pontificia Universidad Javeriana\par
    Facultad de Ingeniería\par
    Trabajo de Grado\par
    \vspace{16pt}
    Bogotá D.C.\par
    \kpFecha\par
    Versión \kpVersion\par
  \end{titlepage}%
  \pagenumbering{roman}%
}

% Cada fila se agrega con \Cambio{versión}{fecha}{descripción}{autor(es)}.
\newenvironment{HistorialDeCambios}{%
  \section*{Historial de Cambios}%
  \renewcommand{\arraystretch}{1.3}%
  \tabularx{\textwidth}{|c|c|X|>{\raggedright\arraybackslash}p{4.2cm}|}
  \hline
  \textbf{Versión} & \textbf{Fecha} & \textbf{Descripción del cambio} & \textbf{Autor(es)} \\
  \hline
}{%
  \endtabularx
  \clearpage
}
\newcommand{\Cambio}[4]{#1 & #2 & #3 & #4 \\ \hline}

% La plantilla PTP pide que la numeración arábiga empiece en la Introducción y que las
% listas de tablas y figuras no entren en la tabla de contenidos.
\newcommand{\Indices}{%
  \tableofcontents
  \clearpage
  \listoftables
  \listoffigures
  \clearpage
  \pagenumbering{arabic}%
}
.

\usepackage[T1]{fontenc}
\usepackage[utf8]{inputenc}
\usepackage[spanish,es-tabla,es-noshorthands,es-lcroman]{babel}
\usepackage{tgtermes}
\usepackage{tgheros}
\usepackage[letterpaper,margin=1in]{geometry}
\usepackage{graphicx}
\usepackage{array}
\usepackage{tabularx}
\usepackage{setspace}
\usepackage{titlesec}
\usepackage{titletoc}
\usepackage{caption}
\usepackage[hidelinks]{hyperref}

\setstretch{1.15}
\setlength{\parindent}{0pt}
\setlength{\parskip}{6pt plus 2pt}

\addto\captionsspanish{%
  \renewcommand{\contentsname}{Tabla de Contenidos}%
  \renewcommand{\listtablename}{Lista de Tablas}%
  \renewcommand{\listfigurename}{Lista de Figuras}%
}

\titleformat{\section}{\normalfont\bfseries\fontsize{12}{14}\selectfont}{\thesection.}{0.3em}{}
\titleformat{\subsection}{\normalfont\bfseries\fontsize{11}{13}\selectfont}{\thesubsection}{0.3em}{}
\titleformat{\subsubsection}{\normalfont\bfseries\normalsize}{\thesubsubsection}{0.3em}{}
\titlespacing*{\section}{0pt}{14pt}{8pt}
\titlespacing*{\subsection}{0pt}{12pt}{6pt}
\titlespacing*{\subsubsection}{0pt}{10pt}{4pt}

\titlecontents{section}[0pt]{\addvspace{2pt}\bfseries}
  {\thecontentslabel.\ }{}{\titlerule*[0.3em]{.}\contentspage}
\titlecontents{subsection}[1.5em]{}
  {\thecontentslabel\ }{}{\titlerule*[0.3em]{.}\contentspage}
\titlecontents{subsubsection}[3em]{}
  {\thecontentslabel\ }{}{\titlerule*[0.3em]{.}\contentspage}
\titlecontents{table}[0.75in]{}
  {\tablename~\thecontentslabel.\ }{}{\titlerule*[0.3em]{.}\contentspage}
\titlecontents{figure}[0.75in]{}
  {\figurename~\thecontentslabel.\ }{}{\titlerule*[0.3em]{.}\contentspage}

\captionsetup{labelsep=period,font=it,justification=centering}

\newcommand{\kpNombre}{}
\newcommand{\kpVersion}{}
\newcommand{\kpFecha}{}
\newcommand{\NombreDocumento}[1]{\renewcommand{\kpNombre}{#1}}
\newcommand{\VersionDocumento}[1]{\renewcommand{\kpVersion}{#1}}
\newcommand{\FechaDocumento}[1]{\renewcommand{\kpFecha}{#1}}

\newcommand{\Portada}{%
  \hypersetup{
    pdftitle={\kpNombre\ - Kit Pagos Colombia},
    pdfauthor={Andres Henao Niño, Joan Emmanuel Orduz Chia, Joshua Prieto Zambrano, David Estevan Rodriguez Jurado}
  }%
  \begin{titlepage}
    \centering
    \setstretch{1}%
    \setlength{\parskip}{0pt}%
    \fontsize{12}{16}\selectfont
    \includegraphics[height=6.96cm]{../comun/logo-javeriana}\par
    \vspace{46pt}
    {\fontsize{14}{16}\selectfont\bfseries\kpNombre\par}
    \vspace{16pt}
    Presentan\par
    \vspace{16pt}
    Andres Henao Niño\par
    Joan Emmanuel Orduz Chia\par
    Joshua Prieto Zambrano\par
    David Estevan Rodriguez Jurado\par
    \vspace{31pt}
    Director\par
    Jaime Andrés Pavlich Mariscal\par
    \vspace{21pt}
    Pontificia Universidad Javeriana\par
    Facultad de Ingeniería\par
    Trabajo de Grado\par
    \vspace{16pt}
    Bogotá D.C.\par
    \kpFecha\par
    Versión \kpVersion\par
  \end{titlepage}%
  \pagenumbering{roman}%
}

% Cada fila se agrega con \Cambio{versión}{fecha}{descripción}{autor(es)}.
\newenvironment{HistorialDeCambios}{%
  \section*{Historial de Cambios}%
  \renewcommand{\arraystretch}{1.3}%
  \tabularx{\textwidth}{|c|c|X|>{\raggedright\arraybackslash}p{4.2cm}|}
  \hline
  \textbf{Versión} & \textbf{Fecha} & \textbf{Descripción del cambio} & \textbf{Autor(es)} \\
  \hline
}{%
  \endtabularx
  \clearpage
}
\newcommand{\Cambio}[4]{#1 & #2 & #3 & #4 \\ \hline}

% La plantilla PTP pide que la numeración arábiga empiece en la Introducción y que las
% listas de tablas y figuras no entren en la tabla de contenidos.
\newcommand{\Indices}{%
  \tableofcontents
  \clearpage
  \listoftables
  \listoffigures
  \clearpage
  \pagenumbering{arabic}%
}
